% Chapter 3 — Requirements Analysis

\chapter{Requirements Analysis}
\label{ch:req}

\section{Who cares about the database?}
Four groups interact with LearnTrack:

\textbf{Students} want one profile, correct enrollment status, visible progress, and quiz scores they can trust. If two screens show different percentages, the database design failed.

\textbf{Instructors} need to publish ordered lessons, attach quizzes, see who enrolled, and skim engagement (watch time, completion, weak quiz averages).

\textbf{Administrators} need user counts, course counts, and sometimes the ability to deactivate accounts or remove bad data---without breaking foreign keys.

\textbf{Developers and DBMS course evaluators} need a single \texttt{ddl.sql} file that rebuilds the same schema anywhere, with constraints documented so behaviour is predictable.

\section{Functional requirements (mapped to tables)}
Table~\ref{tab:funcreq} links plain-language needs to the tables that satisfy them. Use it during vivas: ``Requirement R5 is activity logging; that is \texttt{activity\_log} plus \texttt{activity\_types}.''

\begin{table}[htbp]
\centering
\caption{What the system must do, and where it lives in the schema}
\label{tab:funcreq}
\small
\begin{tabular}{@{}p{8cm}l@{}}
\toprule
\textbf{Requirement (plain wording)} & \textbf{Main tables} \\
\midrule
R1: Every person has one account and a role (student / instructor / admin). & \texttt{users}; instructors also use \texttt{instructors}. \\
R2: Courses belong to an instructor; we can hide a course without deleting history. & \texttt{courses}, \texttt{instructors} \\
R3: Each course has ordered lessons (video, document, quiz marker). & \texttt{content}, \texttt{content\_types} \\
R4: A student joins a course once; status can change (active, completed, dropped). & \texttt{enrollments}, \texttt{enrollment\_statuses} \\
R5: Store fine-grained events: play, pause, skip, plus seconds watched. & \texttt{activity\_log}, \texttt{activity\_types} \\
R6: Store how far each student got on each lesson (0--100\%). & \texttt{content\_progress} \\
R7: Full quiz engine: questions, options, attempts, per-answer rows. & \texttt{quiz}, \texttt{quiz\_questions}, \texttt{question\_options}, \texttt{quiz\_attempts}, \texttt{quiz\_answers} \\
R8: Dashboards show summaries without recomputing giant joins every time. & Materialized views (\texttt{mv\_*}); controller queries \\
\bottomrule
\end{tabular}
\end{table}

\section{Non-functional requirements (what ``good'' means)}
\begin{itemize}[leftmargin=*]
  \item \textbf{Data integrity.} Foreign keys stop dangling enrollments. Unique constraints stop duplicate enrollments for the same student and course.

  \item \textbf{Fair concurrency.} Two enrollment requests at the same second should not create two rows for the same pair. The database \texttt{UNIQUE (user\_id, course\_id)} handles this even if the application forgets to check.

  \item \textbf{Speed for dashboards.} Instructors expect charts to load in seconds. Pre-joined materialized views trade a bit of staleness for fast reads.

  \item \textbf{Maintainability.} Lookup tables (\texttt{activity\_types}, etc.) mean we rename ``skip'' in one place instead of editing thousands of rows.

  \item \textbf{Reset and demo.} The DDL drops and recreates objects in a safe order so coursework and demos can start from a clean slate (\texttt{ddl.sql} + optional \texttt{seed.sql}).

  \item \textbf{Traceability.} Timestamps (\texttt{created\_at}, \texttt{event\_at}, \texttt{attempt\_date}) let us reconstruct timelines when something looks wrong.

  \item \textbf{Room to grow.} Event tables (\texttt{activity\_log}) can accumulate many rows; indexes on user, content, and time support filtered analytics.
\end{itemize}

\section{Requirements traceability tip}
When writing exams or presentations, pick one user story per requirement: e.g.\ ``Student submits quiz'' touches \texttt{quiz\_attempts}, \texttt{quiz\_answers}, and indirectly \texttt{quiz\_questions}, \texttt{question\_options}. Saying the chain out loud proves you understand the schema.
